% =====================================================================
% SAGE26 Paper I -- Section 6 rewrite (Croton's annotations, pp. 10-12).
%
% His "(!?)" ties three notes into one: the Eq. 15 bracket, "So opposite to
% what Dekel intended?", and "state the intended physics, then how eq 17
% produces it".  Answered as one chain: Eq. 15 (a cooling time) -> DB06's
% cooling/compression ratio -> f_stream -> the behaviour it produces.
%
% MODEL CHANGE: fiducial is now DB06 eq. 39 (ColdStreamCeilingOn = 1).  The
% submitted Eq. 17 was the reciprocal of their eq. 38, the SPHERICAL INFALL
% ratio, not eq. 39, the STREAM ratio; it omitted (f M_*/M)^(2/3).
% Verified against the paper in the project folder, p. 11, eqs 37-41.
%
% 8 equations here vs 6 submitted, so downstream numbering shifts by +2.
% See notes at the end.
% =====================================================================

\section{The gas cycle: intermediate-to-high mass systems}

Massive galaxies, groups and clusters of galaxies are treated similarly to
SAGE16, where we assume the infalling gas undergoes shock heating (when
$\mathrm{M_{vir}} \geq \mathrm{M_{shock}}$) and forms a hot halo with an
isothermal density profile at the virial temperature. This becomes the hot gas
reservoir, seen observationally as the intragroup or intracluster medium. Gas
then cools onto the galaxy at a rate dependent on the cooling time.

\subsection{Hot Halo Cooling and Cold Streams}

As a halo grows and evolves past $\mathrm{M_{shock}}$, it transitions from the
free-fall-limited inflow of the circumgalactic medium regime to a slower hot
halo accretion mode. However, this transition does not instantly cut off the
cold gas supply. Even as the shock-heated atmosphere develops, cold streams can
still penetrate through and deliver gas to the galaxy. At this point, a halo can
experience both hot accretion and cold streams delivering gas to the galactic
disk \citep{Keres2005, DekelBirnboim2006, Dekel2009}.

The radiative cooling of hot halo gas proceeds via the standard SAGE16
prescription:

\begin{equation}
\dot{\mathrm{m}}_{\mathrm{cool}} = \frac{1}{2}
\left( \frac{\mathrm{r_{cool}}}{\mathrm{R_{vir}}} \right)
\left( \frac{\mathrm{m'_{hot}}}{t_{\mathrm{dyn}}} \right).
\label{eq:mcool}
\end{equation}

Here $\mathrm{r_{cool}}$ is the cooling radius, at which the cooling time of the
gas equals the dynamical time $t_{\mathrm{dyn}} =
\mathrm{R_{vir}}/\mathrm{V_{vir}}$, and the cooling time of a parcel of hot gas
is approximated by the ratio of its thermal energy to the cooling rate:

\begin{equation}
t_{\mathrm{cool}} = \frac{3}{2}
\frac{\bar{\mu}\, m_{\mathrm{p}}\, k\, \mathrm{T_{vir}}}
     {\rho_{\mathrm{hot}}\, \Lambda(\mathrm{T_{vir}}, Z)},
\label{eq:tcool}
\end{equation}

where $\rho_{\mathrm{hot}}$ is the hot gas density, assumed to be isothermal.
Equation~\ref{eq:tcool} fixes only the location of $\mathrm{r_{cool}}$; the rate
at which gas reaches the galaxy is set by $t_{\mathrm{dyn}}$, following SAGE16.

Cold streams allow part of the accreting gas to pierce the shock-heated
atmosphere, so not all of the halo's gas is available for radiative cooling. The
effective mass is

\begin{equation}
\mathrm{m'_{hot}} \equiv (1 - f_{\mathrm{stream}})\, \mathrm{m_{hot}},
\label{eq:mhot_eff}
\end{equation}

where $f_{\mathrm{stream}}$ is the fraction of the reservoir carried by streams.

\citet{DekelBirnboim2006} decide whether a stream survives by comparing its
cooling time with its compression time. For spherical infall through a
virialised halo that ratio is
$(\mathrm{M_{vir}}/\mathrm{M_{shock}})^{4/3}$ (their equation 38); inside a
stream it is reduced by the density enhancement
$\rho_{\mathrm{stream}}/\rho_{\mathrm{halo}} \sim
(f\,\mathrm{M_*}/\mathrm{M_{vir}})^{-2/3}$, giving their equation 39,

\begin{equation}
\mathcal{R} \equiv
\left( \frac{t_{\mathrm{cool}}}{t_{\mathrm{comp}}} \right)_{\mathrm{stream}} =
\left( \frac{f\, \mathrm{M_*}(z)}{\mathrm{M_{vir}}} \right)^{2/3}
\left( \frac{\mathrm{M_{vir}}}{\mathrm{M_{shock}}} \right)^{4/3},
\label{eq:dbratio}
\end{equation}

where $\mathrm{M_*}(z)$ is the Press--Schechter clustering mass of each
cosmology and $f \sim 3$. Streams penetrate where $\mathcal{R} < 1$. Rather
than applying a binary cut-off, we model the transition across the vicinity of
the threshold in the same way as the feedback-free burst fraction of
Section~13, taking

\begin{equation}
f_{\mathrm{stream}}(\mathrm{M_{vir}}, z) = S
\left[ \frac{-\log \mathcal{R}}{\Delta \log \mathcal{R}} \right],
\label{eq:fstream}
\end{equation}

where $S(x) = (1 + e^{-x})^{-1}$ is the sigmoid of Equation~39 and
$\Delta \log \mathcal{R} = 0.065$~dex, so $f_{\mathrm{stream}} = 1/2$ at the
threshold and approaches unity where streams penetrate freely. Because the two
channels are weighted by $1 - f_{\mathrm{stream}}$ and $f_{\mathrm{stream}}$, a
finite width is what allows shocked and penetrating gas to coexist in one halo;
a strict step would make them exclusive.

The mass and redshift behaviour then follows from Equation~\ref{eq:dbratio}
rather than being imposed separately. Setting it to unity gives the stream
ceiling of \citet{DekelBirnboim2006}, their equation 40,

\begin{equation}
\mathrm{M_{stream}}(z) \sim \frac{\mathrm{M_{shock}}^2}{f\, \mathrm{M_*}(z)},
\label{eq:mstream_ceiling}
\end{equation}

which closes onto $\mathrm{M_{shock}}$ at their equation 41,
$f\,\mathrm{M_*}(z_{\mathrm{crit}}) = \mathrm{M_{shock}}$, giving
$z_{\mathrm{crit}} = 1.20$ for Millennium and $1.01$ for miniUchuu. No explicit
redshift cut is applied. Below $z_{\mathrm{crit}}$ streams are confined to
$\mathrm{M_{vir}} < \mathrm{M_{shock}}$ and shocked haloes cool radiatively;
above it the ceiling climbs steeply as $\mathrm{M_*}(z)$ falls, reaching
$6.4 \times 10^{12}\,\mathrm{M_\odot}$ at $z = 2$ and
$9.6 \times 10^{13}\,\mathrm{M_\odot}$ at $z = 3$, so cold streams supply
essentially every shocked halo at early times. The two channels are most often
comparable near $z \simeq 1.2$, where $\mathrm{M_{stream}}$ sweeps through the
halo mass function.

The hot gas reservoir is thus partitioned into gas that cools radiatively using
the traditional SAGE16 prescription, and cold streams, which cool on the faster
dynamical timescale:

\begin{equation}
\dot{\mathrm{m}}_{\mathrm{stream}} =
\frac{f_{\mathrm{stream}}\, \mathrm{m_{hot}}}{t_{\mathrm{dyn}}}.
\label{eq:mstream}
\end{equation}

The total cooling rate onto the galaxy is thus:

\begin{equation}
\dot{\mathrm{m}}_{\mathrm{total}} =
\dot{\mathrm{m}}_{\mathrm{cool}} + \dot{\mathrm{m}}_{\mathrm{stream}}.
\label{eq:mtotal}
\end{equation}

When $f_{\mathrm{stream}} = 0$ we recover the original SAGE16 cooling
prescription. Below $\mathrm{M_{shock}}$ there is no hot atmosphere to
penetrate, and the inflow model of Section~3.1 applies instead.


% =====================================================================
% NOTES FOR MB -- not for the manuscript
% =====================================================================
%
% 1. WHAT WAS WRONG. Submitted Eq. 17, (Mvir/Mshock)^-4/3 (1+z)/2, is the
%    reciprocal of DB06 eq. 38 -- SPHERICAL INFALL -- with a (1+z)/2 proxy.
%    The stream density enhancement (f M_*/M)^(2/3) is missing, which is the
%    factor that makes a stream a stream. Streams then shut off near 10^12
%    Msun at every epoch instead of tracking M_*(z): at Mvir = 10 Mshock,
%    f_stream was 0.09 at z = 3 where DB06 give 1. That is his "So opposite
%    to what Dekel intended?"; the paper says so explicitly on p. 11:
%    "At high z ... cold streams appear even in M > Mshock haloes ... as
%    long as M < Mstream."
%
% 2. z_crit IS NO LONGER A PARAMETER. It emerges from eq. 41 (1.201
%    Millennium / 1.006 Uchuu, verified against the code's M_*(z) table).
%    Z_CRIT_DB06 = 1.5 survives only for modes 0 and 2.
%
% 3. SMOOTHING is the only addition to DB06. StreamThresholdWidthDex = 0.15;
%    setting 0 gives their step exactly. It moves stellar mass by <1.2% and
%    the SMF by less than Poisson noise, so it is not tuning the results --
%    but with a step, measured coexistence is exactly ZERO at every
%    redshift, which would contradict this subsection's first paragraph.
%    With the sigmoid, the fraction of shocked haloes where each channel
%    supplies >10% of the inflow peaks at ~50% near z = 1.2 and is <1% by
%    z = 0.6 and z = 2.1.  (That measurement is the basis for the single
%    sentence about z ~ 1.2 in the text; the rest stays out of the paper.)
%
% 4. CALIBRATION. Total stellar mass rises 4.9-14.6% (peak near z ~ 2) and
%    the SMF moves up to 0.22 dex at z = 2.1, measured mode 0 -> mode 1 on
%    the full mini-Millennium box with the step; the sigmoid shifts this by
%    <1.2%. Re-measure before quoting. Figure 10 must be regenerated.
%
% 5. KNOWN INCONSISTENCY. r_cool is located from rho0 = HotGas/(4 pi Rvir),
%    the FULL hot gas (model_cooling_heating.c:204), while the mass that
%    cools is m'_hot. For fractional f_stream the cooling flow therefore
%    uses a slightly too-large r_cool. Bites only in the mixed z ~ 1-2
%    window, at the level of the rcool/Rvir prefactor. Fix would be to pass
%    m'_hot into rho0.
%
% 6. Fixed separately: mdot_cool/mdot_stream are now reset each call
%    (model_cooling_heating.c). They were stale whenever a galaxy took an
%    early return or flipped regime -- 1006 galaxies showed phantom streams
%    at z = 1.08. Diagnostic only, but it is a saved output field.
%
% 7. SIGMOID: 0.065 dex IN THE TEXT vs 0.15 dex IN THE CODE -- not an error.
%    The code writes f_stream = 1/(1 + 10^(logR/0.15)), a BASE-10 logistic, so
%    "0.15 dex" there is a literal dex width. Section 13 / Li+24 eq. 3 define
%    the house sigmoid S(x) = (1+e^-x)^-1, base e. Re-expressing the code's
%    form in that S gives exactly S(-logR/0.065), since 0.15/ln(10) = 0.0651.
%    Verified identical to machine precision over log R = -0.3 .. +0.3.
%    The text therefore matches Section 13's notation with NO model change.
%    If you would rather quote 0.15 dex in both sections, the code must switch
%    to exp() -- but that widens the stream transition by ln(10) = 2.3x, from
%    a 0.29 dex to a 0.66 dex 10-90% span, which IS a model change and would
%    need recalibrating. Note the same ln(10) applies to Section 13 and to the
%    regime sigmoid (model_regimes.c:133, exp with 0.1 dex): both are ~2.3x
%    wider than their stated dex. Li+24 call their 0.15 dex "quite arbitrary",
%    so this is a presentation issue, not a physics one -- but a reader WILL
%    assume two "0.15 dex" sigmoids are the same width.
%
% 8. Citation keys are guesses -- check against the .bib.
